%!TEX root = ../../main.tex
\section{초기 우세에 대한 의존: 균형 구간 $\Bforty$}
\label{sec:val-b40}

M-RQ4의 첫 절은 예측이 초기 경기 우세에 얼마나 의존하는가이다. $\PTflex$ 기준선이 이미 $\ppre$의 유연한 함수이므로 주 대비 자체가 초기 우세를 넘는 이득을 잰다. 여기서는 그 이득이 초기 우세가 거의 없는 $\Bforty$ 안에서도 남는지를 각 코호트 안에서 본다(표 \ref{tab:b40}).

%% table 객체: tab:b40 — 출처 RR12 json TEST_B40, bootstrap.B40_q_minus_PT_flex, H_B40_minus_outside; SS table2a_identity, _S_qS_id json
\begin{table}[htbp]
  \centering
  \caption{균형 구간 $\Bforty$ ($0.40 \le \ppre \le 0.60$) 안의 코호트별 결과. 각 블록은 코호트 안에서만 읽는다.}
  \label{tab:b40}
  \small
  \begin{tabular}{@{}llrrrr@{}}
    \toprule
    코호트 & 모델 & $n$ & 경기 & Brier & AUC \\
    \midrule
    T & $q$ & \numTBforty{} & \numTBfortyMatches{} & \qBfortyBrier{} & \qBfortyAUC{} \\
    T & $\PTflex$ & \numTBforty{} & \numTBfortyMatches{} & \ptflexBfortyBrier{} & \ptflexBfortyAUC{} \\
    \multicolumn{6}{@{}l}{\quad $\dBrier$ = \dBrierTBforty{} \ci{\dBrierTBfortyLo}{\dBrierTBfortyHi};\quad $H$ = \hetT{} \ci{\hetTLo}{\hetTHi}} \\
    \midrule
    S & $\qS$ & \numSBforty{} & \numSBfortyMatches{} & \qSBfortyBrier{} & \qSBfortyAUC{} \\
    S & $\PTflexS$ & \numSBforty{} & \numSBfortyMatches{} & \ptflexSBfortyBrier{} & \ptflexSBfortyAUC{} \\
    \multicolumn{6}{@{}l}{\quad $\dBrier$ = \dBrierSBforty{} \ci{\dBrierSBfortyLo}{\dBrierSBfortyHi};\quad $H_S$ = \hetS{} \ci{\hetSLo}{\hetSHi}} \\
    \bottomrule
  \end{tabular}
\end{table}


\paragraph{T.} $\Bforty$(\numTBforty{} 교전)에서 $\dBrier(q - \PTflex) = \dBrierTBforty$ \ci{\dBrierTBfortyLo}{\dBrierTBfortyHi}이다. 구간의 상단이 0 근처에 있으므로 이는 해석 눈금 $\tauMeaningful$보다 작은 \emph{작은 탐색적} 개선이다. 이질성 $H = D_{\Bforty} - D_{\text{outside}}$($D$는 $q$와 $\PTflex$의 행별 제곱오차 차이의 평균)는 \hetT{} \ci{\hetTLo}{\hetTHi}로 구간이 0을 포함하며, 균형 구간 안팎의 이득 차이는 이 표본에서 불확실하다. RR4가 다른 시드로 같은 대비를 다시 추출하면 상단이 $-0.0001$로 인쇄되고 점추정은 같다.

\paragraph{S.} $\Bforty$(\numSBforty{} 교전)에서 $\dBrier(\qS - \PTflexS) = \dBrierSBforty$ \ci{\dBrierSBfortyLo}{\dBrierSBfortyHi}이고 $H_S = \hetS$ \ci{\hetSLo}{\hetSHi}이다. S 안에서 이득은 $\Bforty$ 안에서 바깥보다 작고, $\Bforty$의 이득 자체는 구간이 0을 제외한다. 두 $\Bforty$ 셀은 코호트에서 차지하는 비율(\shareTBforty\% 대 \shareSBforty\%)과 크기(\numTBforty{} 대 \numSBforty{} 행)가 다르며, 각 셀의 결과는 그 코호트 안에서 읽는다.
